\section{Echtzeit Prezesse}
Um zeitkritische Prozesse zuverlässig zu steuern, wird das Echtzeitbetriebssystem
FreeRTOS verwendet. Dies ist notwendig, da \texttt{lvgl} einen handler benötigt,
der alle 5ms aufgerufen werden soll, um die Benutzeroberfläche zu aktualisieren.
Ebenso soll ESP-NOW periodisch nach einer Bridge suchen, um Informationen 
über das Smarthome zu erhalten. Ein RTOS vereinfacht die Implementierung
dieser Prozesse dramatisch.

    \subsection{Struktur}
    In der \texttt{Setup()} Funktion werden drei Tasks erstellt, diese sind 
    zuständig für: 
    \begin{itemize}
        \item lvgl Handler
        \item Input Logik
        \item ESP-NOW
    \end{itemize}

        \subsubsection{lvgl Handler}
        Der lvgl Handler wird alle 5ms aufgerufen und ist für die Aktualisierung
        der Benutzeroberfläche zuständig. Dies geschieht durch den Aufruf
        der lvgl Funktion \texttt{lv\_task\_handler()}. Anschließend wird der 
        Task für 5ms in den Sleep-Mode versetzt.
        \begin{lstlisting}[style=cppCode]
    void loop() { // should only contain lvgl handler
        lv_task_handler();
        vTaskDelay(5 / portTICK_PERIOD_MS);
    }
            \end{lstlisting}

        \subsubsection{Input Logik}
        Der Input Logik Task ist für die Verarbeitung der Eingaben zuständig.
        Diese werden am Rotary Encoder eingelesen und in lvgl Aktionen umgewandelt.

        \subsubsection{ESP-NOW}
        Der ESP-NOW Task ist nur für die Ursprungsuche der Bridge zuständig.
        Diese wird durchgeführt, indem alle 5 Sekunden ein Broadcast-Paket
        gesendet wird. Wird auf dieses Paket geantwortet, wird die 
        MAC-Adresse der Bridge gespeichert und die Suche permanent beendet.

    \subsection{Race Conditions}
    Da Der ESP32-S3 über zwei Kerne verfügt, kann es vorkommen, dass lvgl-Objekte
    bearbeitet werden, während das User-Interface aktualisiert wird. Dies kann
    zu unvorhersehbaren Fehlern und somit zu einem Absturz des Programms führen.
    Um dies zu verhindern, wird ein Mutex verwendet. Dieser wird sowohl
    im lvgl Handler, als auch in der Input Logik aktiviert, sobald ein lvgl-Objekt
    bearbeitet wird.\\~\\
    Vereinfacht kann der Mutex wie folgt verwendet werden:

    \begin{lstlisting}[style=cppCode]
    xSemaphoreTake(lvglMutex, portMAX_DELAY);
    // modify lvgl object
    xSemaphoreGive(lvglMutex);
    \end{lstlisting}

    Wird diese Methode bei jeder Bearbeitung eines lvgl-Objektes verwendet,
    wird sichergestellt, dass immer nur ein Task auf das lvgl-Objekt zugreift.
